\documentclass[11pt,a4paper]{article}
\usepackage[margin=2.2cm]{geometry}
\usepackage{amsmath,amssymb,xcolor,enumitem,booktabs,graphicx,tabularx}
\usepackage[normalem]{ulem}
\setlength{\parskip}{0.6em}\setlength{\parindent}{0pt}
\definecolor{newc}{HTML}{1f6f43}
\newcommand{\new}[1]{\textcolor{newc}{#1}}
\newcommand{\old}[1]{\textcolor{gray}{\sout{#1}}}
\newcommand{\note}[1]{{\small\itshape [#1]}}
\newcommand{\lab}[1]{\textsf{\footnotesize #1}}
\renewcommand{\ref}[1]{\lab{#1}}
\renewcommand{\eqref}[1]{(\lab{#1})}
\newcommand{\nref}[1]{the (\lab{#1})}
\newcommand{\citep}[1]{[\lab{#1}]}
\newcommand{\citet}[1]{\lab{#1}}
\title{Proposed change, ``What an agent is rewarded for''}
\date{24 November 2026, version 1}
\begin{document}\sloppy
\maketitle

\section*{How to read this}
This answers what is left of \texttt{method\_reward\_comments.md} for
\lab{sec:method:rl} in \texttt{4.Implementation.tex}. A4 and the per-task
shaping were done in [A109]. Agreed on 24 November: the B3 arguments, the
rename $n_{\mathrm{blocked}} \to n_{\mathrm{override}}$, dropping
$\omega_i(t)$, the numbers moved down to the step costs (D1), and a table of
the reward terms (E). The rest (A5, B1, B2, B4, C4, D2, D3) is fixed as the
comments say. No value, no term and no behaviour of the reward changes. The
code already calls the weight \texttt{override\_penalty}.

Black text stays, \new{green} is new and \old{struck} text goes.
\note{Notes in brackets} say why. Nothing below is in the thesis yet.

\section*{Part 1. What was checked}
\begin{itemize}[nosep]
  \item van Knippenberg et al.\ (2021), Table~1: move $-0.1$, wait off goal
    $-0.2$, collision $-2$, wait on goal $+0.1$, complete episode $+10$. They
    have no shaping. A collision is possible there.
  \item \texttt{multi\_agent\_env.py:599}: the crowding charge uses
    $\delta_i(t)$ ``at the step the action is chosen in'' (A5).
  \item \lab{eq:congestion}: $\delta_i(t)$ is the share of the field of view
    where another agent stands, between 0 and 1.
  \item The worked example: with $\Phi = -$distance left and $\gamma = 0.99$,
    a step closer gives $\gamma(-(d-1)) + d - 0.1 \approx +0.9$, a wait
    $-0.2$, a step away $\approx -1.1$, each plus $(1-\gamma)$ times the
    distance left. Over a task the shaping sums to about the start distance.
  \item \texttt{rl\_review\_2026-11-10}: override penalty 0.1 and 0 at fleet 32
    had no effect on the collapse. That goes to Appendix~I (Part~4).
\end{itemize}

\section*{Part 2. The section as it would read}

The reward pays an agent for progress through its task and charges it for
time lost. \old{When little of the task is left, a step that brings it closer
to done is worth about $+0.9$, a wait $-0.2$ and a step away $-1.1$. So moving
on is best, and on a single step waiting is cheaper than stepping away.}
\new{A step closer to done pays, a wait costs a little and a step away costs
most.} An agent also earns a bonus for finishing a task, pays extra when
conflict resolution overrides its move, and pays a little \old{for every step
in a crowded neighbourhood} \new{when its neighbourhood is crowded.
Table~\ref{tab:rewardterms} lists the terms, their values and where each
comes from}. \note{D1. The numbers move to the step costs, where $\gamma$ and
$c_{\mathrm{move}}$ are known.}

\new{\textbf{Table, caption:} The terms of the reward.}

\begin{center}\small
\begin{tabularx}{\linewidth}{l l l X}
\toprule
term & weight & value & source \\
\midrule
progress (shaping) & $\gamma$ & 0.99 & \citet{ng1999policy}, $\gamma$ in Appendix~\ref{app:cal:discount} \\
finishing a task & $n_{\mathrm{deliver}}$ & 10 & \citet{vanKnippenberg2021}, paid per task, not per episode \\
being overridden & $n_{\mathrm{override}}$ & 0.5 & the thesis's own, Appendix~\ref{app:cal:reward} \\
a move & $c_{\mathrm{move}}$ & 0.1 & \citet{vanKnippenberg2021}, Table~1 \\
standing still with a task & $c_{\mathrm{wait}}$ & 0.2 & \citet{vanKnippenberg2021}, Table~1 \\
crowding & $n_{\mathrm{crowd}}$ & 0.1 & the thesis's own, Appendix~\ref{app:cal:reward} \\
\bottomrule
\end{tabularx}
\end{center}
\note{E. Values from \lab{tab:rlparams} (Appendix~B), unchanged.}

\new{Progress is counted at every step, because the delivery bonus comes many
steps later and the target is usually outside the field of view. Without it
an agent would learn of its progress only when a task ends. The reward of
\citet{vanKnippenberg2021} has no such term.} Progress is measured by the
distance left in the whole task. \emph{[Rest of the paragraph and
\lab{eq:taskpotential} unchanged.]} \note{B1.}

and the reward is
\begin{equation*}
  \begin{aligned}
  R_i(t) = {}&\gamma\Phi_i(t+1) - \Phi_i(t)
           \;+\; n_{\mathrm{deliver}}\,\mathbb{1}[a_i \text{ completes a task at } t]
           \\
         &\;-\; n_{\mathrm{\old{blocked}\new{override}}}\,\mathbb{1}[a_i \text{ is overridden at } t]
           \;-\; c_{\mathrm{move}}\,\mathbb{1}[a_i \text{ moves at } t]
           \\
         &\;-\; c_{\mathrm{wait}}\,\old{\omega_i(t)}\new{\mathbb{1}[a_i \text{ serves a task and stands still at } t]}
           \;-\; n_{\mathrm{crowd}}\,\delta_i(t) .
  \end{aligned}
  \tag{\lab{eq:reward}}
\end{equation*}
\note{C1 and C3. Same terms and values, two names changed.}

\emph{[The per-task paragraph, ``The reason is that an agent avoids a
charge\dots'' and the paragraph on potential-based shaping are unchanged.]}

Standing still costs time, and the reward says so. An agent that serves a
task and stands still pays the wait cost $c_{\mathrm{wait}}$\old{ of
\nref{eq:onestepcost}}, whether it chose to wait or was \old{stopped}
\new{overridden}. \old{Formally, $\omega_i(t) = \omega_j(t)$ of
\nref{eq:split} for the task $\tau_j$ with $\sigma(j) = i$, and $\omega_i(t)
= 0$ when the agent serves no task, so an idle agent pays nothing.} \new{It
is the same standing still that the blocked indicator $\omega_j(t)$ of the
service-time split~\eqref{eq:split} counts for the task, and an idle agent
pays nothing.} Being \old{stopped} \new{overridden} therefore costs
$c_{\mathrm{wait}} + n_{\mathrm{override}}$ and a chosen wait costs
$c_{\mathrm{wait}}$ alone, so \old{yielding} \new{a chosen wait} is never
dearer than being \old{stopped} \new{overridden}. \note{C1, C3, C4. One word,
``overridden'', everywhere.}

\new{There is no collision penalty, unlike the $-2$ of
\citet{vanKnippenberg2021}. There a collision happens and ends the episode,
so the penalty is what keeps agents apart. Here conflict resolution turns
every clashing move into a wait, so no collision ever happens. Without a
penalty, though, a refused move would cost exactly what a chosen wait costs,
and the reward would never show the agent that it asked for a vertex that
was taken. The override term puts that cost on the agent whose move is
refused. It is the thesis's own. It only has to make an override dearer
than a chosen wait, which any positive weight does, and its size is set in
Appendix~\ref{app:cal:reward}.} A centralised agent never pays it, because
the claim mask keeps its moves apart. A section-based agent pays it only at
the border between two zones, and a decentralised agent at any clash.
\note{B3, as agreed. That the same ordering favours waiting in front of a
workstation is said in ``When moves still clash''.}

Moving costs a little too. Every move pays $c_{\mathrm{move}}$, with or
without a task, so an agent with no reason to move stays put. The two step
costs are those of \citet{vanKnippenberg2021}, a move $0.1$ and a wait
$0.2$, which they set from common practice and which are taken as they are.
\old{Their reward has no shaping term.} \old{Here the shaping term pays about
one for each step that shortens the distance left in the task and charges
about one for each step that lengthens it, which gives the $+0.9$, $-0.2$
and $-1.1$ above.} \new{With the discount $\gamma = 0.99$, the shaping term
pays about one for a step that shortens the distance left in the task and
charges about one for a step that lengthens it. When little of the task is
left, a step closer is then worth about $+0.9$, a wait $-0.2$ and a step
away $-1.1$.} With more of the task left, every action gets the same small
extra amount, $1-\gamma$ times the distance left. So the order of the three
stays the same.

\new{Take a task of 12 steps done without a stop. The shaping pays about 12,
the delivery bonus 10, and the twelve moves cost 1.2, so the task is worth
about 20.8. Every step the agent stands still on the way costs 0.2 more,
and 0.7 if it was overridden.} \note{D3, a worked example on the values of
the table.}

These are the values of one step. Over a whole route the order of waiting
and a detour turns round. \emph{[Rest of the paragraph unchanged.]}

\old{The last term, the crowding term, charges for crowding.} \new{The
crowding term charges an agent in proportion to how full its field of view
is, $\delta_i(t)$~\eqref{eq:congestion}. The network already sees
$\delta_i(t)$, but seeing a crowd does not make it costly, and the term
does. It is the thesis's own. It is read at the step the move is chosen in,
so a move changes only the next step's charge.} Even within a task it can
change which behaviour is best, and that is accepted because it gives the
policy a direct signal to avoid crowds. The delivery bonus, the override
\old{cost} \new{weight} and the crowding weight are set in
Appendix~\ref{app:cal:reward}. \note{A5, B4, C4.}

\section*{Part 3. Notation}
\begin{itemize}[nosep]
  \item Both notation tables: $n_{\mathrm{deliver}},
    \old{n_{\mathrm{blocked}}}\new{n_{\mathrm{override}}}, n_{\mathrm{crowd}}$.
  \item \texttt{A.Notation.tex}: the blocked indicator row drops
    ``$\omega_i(t)$'' and keeps $\omega_j(t)$.
  \item \texttt{B.Parameters.tex} (\lab{tab:rlparams}) and
    \texttt{I.Reward.tex}: $n_{\mathrm{blocked}} \to n_{\mathrm{override}}$.
\end{itemize}

\section*{Part 4. Appendix I}
After ``The override penalty only has to make an override cost more than a
chosen wait, which any $n_{\mathrm{override}} > 0$ does, so its size is
open.'':

\new{Two training runs at 32 agents set it to 0.1 and to 0, and neither
changed the collapse those runs were testing. So the penalty did not cause
the collapse. Whether it helps was not tested.} \note{rl\_review\_2026-11-10,
README rows \texttt{f32\_override\_0.1} and \texttt{f32\_override\_0.0}.}

\section*{Questions}
None. The choices were made in chat on 24 November.

\end{document}
